\documentclass[aspectratio=169]{beamer}

\usepackage[utf8]{inputenc}
\usepackage[T1]{fontenc}
\usepackage{lmodern}
\usepackage{booktabs}
\usepackage{graphicx}
\usepackage{xcolor}
\usepackage{hyperref}

\usetheme{Madrid}
\usecolortheme{seagull}

\title{Vers une approche adaptative combinant IA générative et gamification}
\subtitle{Apprentissage immersif et personnalisé dans les environnements e-learning}
\author{Zouhair El Hadiq}
\institute{Professeur agrégé d'informatique -- CPGE Bab Sahra (Guelmim)\\Fondateur de la plateforme \texttt{cpgeacademy.org}}
\date{Audition doctorale avec Pr. Zahi Jarir -- 2025}

\begin{document}

\begin{frame}
  \titlepage
\end{frame}

\begin{frame}{Pourquoi ce projet ?}
  \begin{itemize}
    \item Beaucoup d'apprenants décrochent encore dans les cours en ligne.
    \item Les profils des étudiants CPGE sont très variés : un même contenu ne suffit plus.
    \item L'IA générative permet de créer des explications et exercices à la demande.
    \item La gamification aide à garder le rythme sur la durée.
  \end{itemize}
\end{frame}

\begin{frame}{Qui je suis}
  \begin{itemize}
    \item Professeur agrégé d'informatique à la CPGE Bab Sahra (Guelmim).
    \item Ingénieur en modélisation scientifique (EMI Rabat).
    \item Fondateur de \texttt{cpgeacademy.org}, plateforme e-learning pour les CPGE.
    \item Habitué à accompagner des étudiants vers les concours et à suivre leurs progrès.
  \end{itemize}
\end{frame}

\begin{frame}{Sujet proposé}
  \begin{block}{Question centrale}
    Comment offrir à chaque étudiant un parcours en ligne qui s'adapte à son niveau et à sa motivation, grâce à l'IA générative et à la gamification ?
  \end{block}
  \begin{itemize}
    \item Le projet est porté par le Pr. Jarir et le laboratoire d'Informatique et de Systèmes Intelligents.
    \item Mon apport : un terrain réel (CPGE), une connaissance fine des besoins et une plateforme prête à évoluer.
  \end{itemize}
\end{frame}

\begin{frame}{Vision simplifiée}
  \begin{itemize}
    \item \textbf{Observer} : analyser les réponses et le rythme de chaque étudiant.
    \item \textbf{Adapter} : proposer des activités générées par l'IA (explications, quiz, cas pratiques) au bon niveau.
    \item \textbf{Motiver} : fixer des défis, badges et objectifs adaptés au profil.
    \item \textbf{Accompagner} : donner aux enseignants un tableau de bord clair pour suivre la progression.
  \end{itemize}
\end{frame}

\begin{frame}{Plan de travail sur trois ans}
  \begin{itemize}
    \item \textbf{Année 1} : revue des pratiques, co-construction du modèle pédagogique, premiers tests sur cpgeacademy.
    \item \textbf{Année 2} : développement du prototype, intégration de l'IA générative, scénarios gamifiés.
    \item \textbf{Année 3} : expérimentations élargies (CPGE et formation continue), mesure des effets et publications.
  \end{itemize}
\end{frame}

\begin{frame}{Bénéfices attendus}
  \begin{itemize}
    \item Pour le laboratoire : un projet ancré sur le terrain, des données réelles et des retombées visibles.
    \item Pour les étudiants : plus de motivation, des feedbacks rapides et un parcours personnalisé.
    \item Pour les partenaires : démonstrateurs prêts à être adaptés à d'autres contextes (langues, formation pro).
  \end{itemize}
\end{frame}

\begin{frame}{Conditions de réussite}
  \begin{itemize}
    \item Encadrement régulier avec le Pr. Jarir pour valider chaque étape.
    \item Comité pédagogique CPGE pour tester et ajuster la solution.
    \item Charte d'usage de l'IA et anonymisation des données pour rassurer étudiants et familles.
    \item Accès aux ressources techniques du laboratoire pour héberger et sécuriser le prototype.
  \end{itemize}
\end{frame}

\begin{frame}{Conclusion}
  \begin{itemize}
    \item Le sujet correspond aux besoins actuels de nos étudiants et aux priorités du laboratoire.
    \item Mon profil combine pratique de terrain, maîtrise technique et envie de partager les résultats.
    \item Je suis prêt à lancer un premier atelier laboratoire--CPGE dès la rentrée pour co-construire le prototype.
  \end{itemize}
\end{frame}

\end{document}
